\section{Related Work}
\label{sec:related_work}

\subsection{Industrial Anomaly Detection and Localization}
Unsupervised visual anomaly detection on industrial benchmarks such as MVTec AD \cite{bergmann2019mvtec, czimmermann2020visual} has predominantly evolved along reconstruction-based and feature-embedding paradigms. Reconstruction-based models train autoencoders or generative networks exclusively on nominal images and localize defects via pixel-wise residuals $\|I - \hat{I}\|$, yet often suffer from boundary blurring and over-generalized reconstruction of salient flaws. Feature-embedding approaches mitigate blur by modeling frozen multi-scale descriptors: PaDiM \cite{defard2021padim} fits per-patch multivariate Gaussians, while PatchCore \cite{roth2022towards} builds coreset-subsampled memory banks for nearest-neighbor scoring. More recently, SimpleNet \cite{liu2023simplenet} and EfficientAD \cite{batzner2024efficientad} have substantially accelerated inference toward millisecond-level latencies via feature-space discriminators and lightweight student--teacher distillation. Nevertheless, even these fast anomaly detectors still output coarse continuous anomaly heatmaps rather than crisp polygonal instance masks needed for robotic rejection or defect metrology, and they lack an explicit training-time constraint to suppress false alarms on benign nominal variations.

Conversely, single-stage instance segmentation architectures such as the YOLO family \cite{ultralytics2024yolo11} are widely favored for inline industrial inspection because they directly predict vector-ready instance masks at high throughput. However, supervised YOLO-based defect segmenters depend on extensive pixel-annotated defect corpora, limiting their applicability when only nominal training samples are available.

\subsection{Self-Supervised Anomaly Synthesis}
Bridging unsupervised anomaly detection with supervised segmentation has motivated self-supervised anomaly synthesis from nominal images alone. CutPaste \cite{li2021cutpaste} pastes rotated rectangular or scar-like patches onto source images, whereas DRAEM \cite{zavrtanik2021draem} generates organic shapes via multi-octave Perlin noise \cite{perlin1985image} with linear opacity blending, and Natural Synthetic Anomalies (NSA) \cite{schluter2022natural} employs Poisson image editing \cite{perez2003poisson} to eliminate artificial boundary seams. Recent advances have further refined synthetic supervision: DeSTSeg \cite{zhang2023destseg} couples synthetic defect corruption with a segmentation-guided denoising student--teacher network, while RealNet \cite{zhang2024realnet} synthesizes realistic, strength-controllable anomalies via diffusion models paired with anomaly-aware feature selection.

Despite strong progress, existing spatial anomaly synthesizers share two fundamental scientific limitations:
\begin{enumerate}
    \item \textbf{Purely Spatial Operation:} Manipulating images strictly as spatial pixel grids leaves synthetic defects vulnerable to masking by high-energy periodic background lattices (e.g., woven fabrics or metallic wire meshes).
    \item \textbf{False-Alarm Amplification:} Training discriminative segmentation heads exclusively on synthetically perturbed images induces over-sensitivity to normal surface grain and lighting shifts unless an explicit negative-sample preservation invariant is enforced.
\end{enumerate}

\subsection{Frequency-Domain Inspection and Optical Robustness}
Classical 2D Discrete Fourier (2D-DFT) and Discrete Wavelet Transforms (2D-DWT) \cite{gonzalez2018digital, mallat2008wavelet} offer interpretable decompositions of surface texture, isolating periodic background lattices into conjugate-symmetric spectral impulses and separating multi-resolution edge sub-bands ($LH, HL, HH$). Recent industrial anomaly detection studies have begun integrating frequency cues into deep networks: FAIR \cite{liu2025fair} employs high-pass frequency filtering to guide image restoration and prevent identity-shortcut reconstruction, while Liang \textit{et al.} \cite{liang2023omnifrequency} introduce omni-frequency channel-selection representations within deep feature extractors. However, these methods embed frequency processing inside heavy reconstruction decoders or deep feature-selection networks, whereas our approach applies coordinate-preserving pixel-level frequency preprocessing directly ahead of spatial synthesis to train a lightweight single-stage polygonal segmenter.

Concurrently, robustness under non-ideal optical conditions has become critical for factory deployment. Corruption-robustness benchmarks---pioneered by ImageNet-C \cite{hendrycks2019robustness} and extended to industrial surfaces in MVTec-AD-C \cite{chen2023mvtecadc}---reveal substantial sensitivity of spatial detectors to sensor noise, motion blur, and lens defocus. Coupling spatial synthesis with Haar 2D-DWT detail boosting provides a principled frequency-domain defense against low-pass optical blur.

\textbf{Positioning:} Unlike prior methods that output coarse anomaly heatmaps, embed frequency cues only inside heavy reconstruction backbones, or synthesize defects purely in the spatial domain, our framework unifies coordinate-preserving pixel-level frequency filtering (2D-FFT notch and 2D-DWT modulation) with foreground-clamped spatial synthesis, explicit negative-sample false-alarm suppression, and single-stage polygonal instance segmentation via YOLOv11-seg.
